\documentclass{article}
\usepackage{amsmath,amssymb,amsthm}
\usepackage{algorithm,algorithmic}
\usepackage{enumitem}
\usepackage{hyperref}
\newtheorem{theorem}{Theorem}
\newtheorem{lemma}{Lemma}
\newtheorem{assumption}{Assumption}
\newtheorem{corollary}{Corollary}
\newcommand{\prox}{\operatorname{prox}}
\newcommand{\grad}{\nabla}
\newcommand{\R}{\mathbb{R}}
\newcommand{\EE}{\mathbb{E}}
\newcommand{\norm}[1]{\left\|#1\right\|}
\newcommand{\inner}[1]{\left\langle#1\right\rangle}
\begin{document}

\section*{Algorithm Name}
\textbf{Proximal Gradient Descent (PGD) for Non‑Convex Composite Optimization}

\section*{Mathematical Formulation}
We consider the composite optimization problem  

\[
\min_{x\in\R^{d}} \; F(x)\;\;,\qquad 
F(x):=f(x)+g(x),
\]

where  

\begin{itemize}[noitemsep]
\item $f:\R^{d}\to\R$ is differentiable and $L$‑smooth (i.e. $\|\grad f(x)-\grad f(y)\|\le L\|x-y\|$ for all $x,y$).  
\item $g:\R^{d}\to\R\cup\{+\infty\}$ is proper, lower‑semicontinuous and \emph{convex} so that its proximal operator  

\[
\prox_{\eta g}(z):=\arg\min_{u\in\R^{d}}\Big\{ g(u)+\frac{1}{2\eta}\|u-z\|^{2}\Big\}
\]

is well defined and \emph{non‑expansive}: $\|\prox_{\eta g}(z)-\prox_{\eta g}(z')\|\le\|z-z'\|$ for all $z,z'$.
\end{itemize}

Given a stepsize $\eta>0$, the PGD iteration is  

\[
\boxed{x_{k+1}= \prox_{\eta g}\bigl(x_{k}-\eta \,\grad f(x_{k})\bigr)}\qquad k=0,1,\dots
\tag{1}
\]

Define the \emph{gradient mapping}  

\[
G_{\eta}(x):=\frac{1}{\eta}\Bigl(x-\prox_{\eta g}\bigl(x-\eta \grad f(x)\bigr)\Bigr).
\tag{2}
\]

When $g\equiv0$, $G_{\eta}(x)=\grad f(x)$. The norm $\|G_{\eta}(x_{k})\|$ measures stationarity of $x_{k}$ for the non‑convex problem.

\section*{Pseudocode}
\begin{algorithm}[H]
\caption{Proximal Gradient Descent}
\begin{algorithmic}[1]
\STATE \textbf{Input:} stepsize $\eta\in(0,1/L]$, initial point $x_{0}\in\R^{d}$, maximal iterations $T$.
\FOR{$k=0$ \TO $T-1$}
    \STATE $z_{k}=x_{k}-\eta\,\grad f(x_{k})$ \hfill \# gradient step on $f$
    \STATE $x_{k+1}= \prox_{\eta g}(z_{k})$ \hfill \# proximal step on $g$
\ENDFOR
\STATE \textbf{Output:} $\{x_{k}\}_{k=0}^{T}$ (or the last iterate $x_{T}$)
\end{algorithmic}
\end{algorithm}

\section*{Dry Run Example}
Consider the scalar problem  

\[
f(w)=(w-3)^{2},\qquad g\equiv0,\qquad \eta=0.1,\qquad w_{0}=0 .
\]

Because $g=0$, $\prox_{\eta g}(z)=z$. The iteration reduces to ordinary gradient descent:
$w_{k+1}=w_{k}-\eta\,\grad f(w_{k})$ with $\grad f(w)=2(w-3)$.

\begin{center}
\begin{tabular}{c|c|c|c}
$k$ & $w_{k}$ & $\grad f(w_{k})$ & $F(w_{k})=(w_{k}-3)^{2}$\\ \hline
0 & $0$ & $2(0-3)=-6$ & $9$\\
1 & $0-0.1(-6)=0.6$ & $2(0.6-3)=-4.8$ & $(0.6-3)^{2}=5.76$\\
2 & $0.6-0.1(-4.8)=1.08$ & $2(1.08-3)=-3.84$ & $(1.08-3)^{2}=3.6864$\\
3 & $1.08-0.1(-3.84)=1.464$ & $2(1.464-3)=-3.072$ & $(1.464-3)^{2}=2.3705$
\end{tabular}
\end{center}
The objective value decreases at each iteration, illustrating the descent property proved below.

\newpage
\section*{Assumptions}
\begin{assumption}[Smoothness of $f$]\label{ass:smooth}
$f$ is differentiable and $L$‑Lipschitz gradient:
\[
\|\grad f(x)-\grad f(y)\|\le L\|x-y\|\quad\forall x,y\in\R^{d}.
\]
\end{assumption}
\begin{assumption}[Convexity and Proximal Regularity of $g$]\label{ass:g}
$g$ is proper, closed, convex and its proximal operator $\prox_{\eta g}$ is single‑valued for every $\eta>0$.
\end{assumption}
\begin{assumption}[Stepsize]\label{ass:step}
The stepsize satisfies $0<\eta\le 1/L$.
\end{assumption}
\begin{assumption}[Bounded Below]\label{ass:bounded}
The composite objective $F$ is bounded below, i.e. $F^{*}:=\inf_{x}F(x)>-\infty$.
\end{assumption}

\section*{Main Convergence Theorem}
\begin{theorem}[Convergence of PGD]\label{thm:main}
Let $\{x_{k}\}$ be generated by \eqref{1} under Assumptions \ref{ass:smooth}–\ref{ass:step}. Then for any $T\ge1$,
\[
\frac{1}{T}\sum_{k=0}^{T-1}\norm{G_{\eta}(x_{k})}^{2}
\;\le\;
\frac{2\bigl(F(x_{0})-F^{*}\bigr)}{\eta\,T}.
\tag{3}
\]
Consequently $\displaystyle\min_{0\le k<T}\norm{G_{\eta}(x_{k})}^{2}=O\!\left(\frac{1}{T}\right)$.
\end{theorem}

\section*{Theoretical Properties}
\begin{enumerate}[noitemsep]
\item \textbf{Stationarity guarantee.} Since $G_{\eta}(x)=0$ iff $0\in\grad f(x)+\partial g(x)$, the bound (3) implies that at least one iterate is an $\varepsilon$‑approximate stationary point with $\varepsilon^{2}=O(1/T)$.
\item \textbf{Stepsize sensitivity.} The rate depends linearly on $\eta$; choosing the largest admissible $\eta=1/L$ yields the smallest constant.
\item \textbf{Comparison to SGD.} Unlike stochastic methods, PGD is deterministic and achieves a deterministic $O(1/T)$ rate for the gradient mapping without variance terms.
\item \textbf{Extension to stochastic gradients.} If $g$ remains convex and the stochastic gradient satisfies unbiasedness and bounded variance, the same proof skeleton leads to an $O(1/\sqrt{T})$ rate after taking expectations.
\end{enumerate}

\section*{Discussion}
The theorem shows that proximal gradient descent enjoys the same $O(1/T)$ stationarity rate as gradient descent for smooth non‑convex functions, while allowing the inclusion of a convex regularizer $g$ handled via the proximal map. The proof crucially uses two ingredients:
\begin{itemize}[noitemsep]
\item the \emph{descent lemma} for $L$‑smooth $f$,
\item the \emph{non‑expansiveness} of $\prox_{\eta g}$ (which holds for convex $g$).
\end{itemize}
Both are invoked in the algebraic manipulations below.

\newpage
\section*{Preliminaries (Definitions and Lemmas)}
\begin{lemma}[Descent Lemma]\label{lem:descent}
If $f$ has $L$‑Lipschitz gradient, then for all $x,y\in\R^{d}$
\[
f(y)\le f(x)+\inner{\grad f(x)}{y-x}+\frac{L}{2}\|y-x\|^{2}.
\tag{4}
\]
\end{lemma}
\begin{proof}
Standard consequence of the integral form of the gradient and the Lipschitz condition. \hfill $\square$
\end{proof}

\begin{lemma}[Optimality Condition of the Proximal Operator]\label{lem:prox-opt}
For any $z\in\R^{d}$ and $\eta>0$, $u=\prox_{\eta g}(z)$ iff
\[
0\in \frac{1}{\eta}(u-z)+\partial g(u)
\quad\Longleftrightarrow\quad
z-u\in \eta\,\partial g(u).
\tag{5}
\]
\end{lemma}
\begin{proof}
The proximal mapping solves a convex optimization problem; the first‑order optimality condition yields \eqref{5}. \hfill $\square$
\end{proof}

\begin{lemma}[Non‑expansiveness]\label{lem:nonexp}
If $g$ is convex, then for any $\eta>0$,
\[
\|\prox_{\eta g}(z)-\prox_{\eta g}(z')\|\le\|z-z'\|\quad\forall\,z,z'\in\R^{d}.
\tag{6}
\]
\end{lemma}
\begin{proof}
Follows from the firm non‑expansiveness of proximal operators of convex functions. \hfill $\square$
\end{proof}

\section*{Main Proof (Step‑by‑Step Derivation)}
We prove Theorem \ref{thm:main}. Throughout we set $x_{k+1}= \prox_{\eta g}(x_{k}-\eta\grad f(x_{k}))$ and denote $z_{k}=x_{k}-\eta\grad f(x_{k})$.

\noindent\textbf{Step 1 (Apply the descent lemma to $f$).}\\
Using Lemma \ref{lem:descent} with $x=x_{k}$ and $y=x_{k+1}$,
\[
f(x_{k+1})\le f(x_{k})+\inner{\grad f(x_{k})}{x_{k+1}-x_{k}}
+\frac{L}{2}\|x_{k+1}-x_{k}\|^{2}.
\tag{7}
\]
\textbf{[Step 1]}

\noindent\textbf{Step 2 (Rewrite the inner product).}\\
Since $z_{k}=x_{k}-\eta\grad f(x_{k})$, we have $\grad f(x_{k})=\frac{1}{\eta}(x_{k}-z_{k})$. Substituting into \eqref{7} gives
\[
f(x_{k+1})\le f(x_{k})+\frac{1}{\eta}\inner{x_{k}-z_{k}}{x_{k+1}-x_{k}}
+\frac{L}{2}\|x_{k+1}-x_{k}\|^{2}.
\tag{8}
\]
\textbf{[Step 2]}

\noindent\textbf{Step 3 (Use the proximal optimality condition).}\\
From Lemma \ref{lem:prox-opt} applied to $z_{k}$,
\[
z_{k}-x_{k+1}\in \eta\,\partial g(x_{k+1}) .
\tag{9}
\]
Thus there exists $v_{k+1}\in\partial g(x_{k+1})$ such that $z_{k}-x_{k+1}= \eta v_{k+1}$, i.e. $x_{k+1}=z_{k}-\eta v_{k+1}$.
\textbf{[Step 3]}

\noindent\textbf{Step 4 (Express $x_{k+1}-x_{k}$).}\\
Using $z_{k}=x_{k}-\eta\grad f(x_{k})$ and \eqref{9},
\[
x_{k+1}-x_{k}= -\eta\grad f(x_{k}) -\eta v_{k+1}.
\tag{10}
\]
\textbf{[Step 4]}

\noindent\textbf{Step 5 (Plug \eqref{10} into \eqref{8}).}\\
The inner product term becomes
\[
\inner{x_{k}-z_{k}}{x_{k+1}-x_{k}}
= \inner{\eta\grad f(x_{k})}{-\eta\grad f(x_{k})-\eta v_{k+1}}
= -\eta^{2}\|\grad f(x_{k})\|^{2}-\eta^{2}\inner{\grad f(x_{k})}{v_{k+1}} .
\]
Hence
\[
f(x_{k+1})\le f(x_{k}) -\eta\|\grad f(x_{k})\|^{2}
-\eta\inner{\grad f(x_{k})}{v_{k+1}}
+\frac{L}{2}\|x_{k+1}-x_{k}\|^{2}.
\tag{11}
\]
\textbf{[Step 5]}

\noindent\textbf{Step 6 (Add $g(x_{k+1})$ and use convexity of $g$).}\\
Convexity of $g$ yields
\[
g(x_{k+1})\le g(x_{k})+\inner{v_{k+1}}{x_{k+1}-x_{k}} .
\tag{12}
\]
Adding \eqref{12} to \eqref{11} gives
\[
F(x_{k+1})\le F(x_{k}) -\eta\|\grad f(x_{k})\|^{2}
-\eta\inner{\grad f(x_{k})}{v_{k+1}}
+\inner{v_{k+1}}{x_{k+1}-x_{k}}
+\frac{L}{2}\|x_{k+1}-x_{k}\|^{2}.
\tag{13}
\]
\textbf{[Step 6]}

\noindent\textbf{Step 7 (Replace $x_{k+1}-x_{k}$ using \eqref{10}).}\\
From \eqref{10}, $x_{k+1}-x_{k}= -\eta\bigl(\grad f(x_{k})+v_{k+1}\bigr)$. Therefore
\[
\inner{v_{k+1}}{x_{k+1}-x_{k}}
= -\eta\inner{v_{k+1}}{\grad f(x_{k})}
-\eta\|v_{k+1}\|^{2}.
\tag{14}
\]
Also,
\[
\|x_{k+1}-x_{k}\|^{2}= \eta^{2}\|\grad f(x_{k})+v_{k+1}\|^{2}
= \eta^{2}\bigl(\|\grad f(x_{k})\|^{2}+2\inner{\grad f(x_{k})}{v_{k+1}}+\|v_{k+1}\|^{2}\bigr).
\tag{15}
\]
Substituting \eqref{14} and \eqref{15} into \eqref{13} yields
\[
\begin{aligned}
F(x_{k+1})
&\le F(x_{k})
-\eta\|\grad f(x_{k})\|^{2}
-\eta\inner{\grad f(x_{k})}{v_{k+1}}  \\
&\quad -\eta\inner{v_{k+1}}{\grad f(x_{k})}
-\eta\|v_{k+1}\|^{2}
+\frac{L}{2}\eta^{2}\bigl(\|\grad f(x_{k})\|^{2}+2\inner{\grad f(x_{k})}{v_{k+1}}+\|v_{k+1}\|^{2}\bigr).
\end{aligned}
\tag{16}
\]
\textbf{[Step 7]}

\noindent\textbf{Step 8 (Collect like terms).}\\
Group the coefficients of $\|\grad f(x_{k})\|^{2}$, $\inner{\grad f(x_{k})}{v_{k+1}}$, and $\|v_{k+1}\|^{2}$:
\[
\begin{aligned}
F(x_{k+1})
&\le F(x_{k})
-\bigl(\eta-\tfrac{L}{2}\eta^{2}\bigr)\|\grad f(x_{k})\|^{2} \\
&\quad -\bigl(2\eta-\!L\eta^{2}\bigr)\inner{\grad f(x_{k})}{v_{k+1}}
-\bigl(\eta-\tfrac{L}{2}\eta^{2}\bigr)\|v_{k+1}\|^{2}.
\end{aligned}
\tag{17}
\]
\textbf{[Step 8]}

\noindent\textbf{Step 9 (Apply the stepsize condition).}\\
Since $0<\eta\le 1/L$, we have $L\eta\le1$ and consequently
\[
\eta-\tfrac{L}{2}\eta^{2}\ge \frac{\eta}{2},\qquad
2\eta-L\eta^{2}\ge \eta .
\tag{18}
\]
Using \eqref{18} in \eqref{17} gives the simplified inequality
\[
F(x_{k+1})\le F(x_{k})-\frac{\eta}{2}\Bigl(\|\grad f(x_{k})\|^{2}+ \|v_{k+1}\|^{2}\Bigr)
-\eta\inner{\grad f(x_{k})}{v_{k+1}} .
\tag{19}
\]
\textbf{[Step 9]}

\noindent\textbf{Step 10 (Express the right‑hand side via the gradient mapping).}\\
Recall the definition \eqref{2}:
\[
G_{\eta}(x_{k})=\frac{1}{\eta}\bigl(x_{k}-x_{k+1}\bigr)
= \grad f(x_{k})+v_{k+1}.
\tag{20}
\]
Hence
\[
\|G_{\eta}(x_{k})\|^{2}= \|\grad f(x_{k})\|^{2}+2\inner{\grad f(x_{k})}{v_{k+1}}+\|v_{k+1}\|^{2}.
\tag{21}
\]
Rearranging \eqref{21},
\[
\|\grad f(x_{k})\|^{2}+ \|v_{k+1}\|^{2}+2\inner{\grad f(x_{k})}{v_{k+1}}
= \|G_{\eta}(x_{k})\|^{2}.
\tag{22}
\]
Insert \eqref{22} into \eqref{19}:
\[
F(x_{k+1})\le F(x_{k})-\frac{\eta}{2}\|G_{\eta}(x_{k})\|^{2}.
\tag{23}
\]
\textbf{[Step 10]}

\noindent\textbf{Step 11 (Telescoping sum).}\\
Sum \eqref{23} for $k=0,\dots,T-1$:
\[
F(x_{T})\le F(x_{0})-\frac{\eta}{2}\sum_{k=0}^{T-1}\|G_{\eta}(x_{k})\|^{2}.
\tag{24}
\]
Since $F(x_{T})\ge F^{*}$ by Assumption \ref{ass:bounded},
\[
\frac{1}{T}\sum_{k=0}^{T-1}\|G_{\eta}(x_{k})\|^{2}
\le \frac{2\bigl(F(x_{0})-F^{*}\bigr)}{\eta\,T}.
\tag{25}
\]
This is exactly the bound \eqref{3} claimed in Theorem \ref{thm:main}.
\textbf{[Step 11]}

The proof is complete. \hfill $\square$

\section*{Rate Derivation (With Constants)}
From \eqref{25} we obtain the explicit constant
\[
C:=\frac{2\bigl(F(x_{0})-F^{*}\bigr)}{\eta},
\qquad
\frac{1}{T}\sum_{k=0}^{T-1}\norm{G_{\eta}(x_{k})}^{2}\le\frac{C}{T}.
\]
Choosing the maximal admissible stepsize $\eta=1/L$ minimizes $C$:
\[
C_{\max}=2L\bigl(F(x_{0})-F^{*}\bigr).
\]
Thus the optimal rate is
\[
\min_{0\le k<T}\norm{G_{1/L}(x_{k})}^{2}
\le \frac{2L\bigl(F(x_{0})-F^{*}\bigr)}{T}=O\!\left(\frac{1}{T}\right).
\]

\section*{Special Cases}
\begin{enumerate}[noitemsep]
\item \textbf{Pure gradient descent ($g\equiv0$).} Then $G_{\eta}(x)=\grad f(x)$ and \eqref{23} reduces to the classic smooth‑non‑convex descent inequality $f(x_{k+1})\le f(x_{k})-\frac{\eta}{2}\|\grad f(x_{k})\|^{2}$, yielding the same $O(1/T)$ stationarity rate.
\item \textbf{Indicator of a convex set $C$: $g=\iota_{C}$.} The proximal step becomes Euclidean projection onto $C$. The theorem therefore covers projected gradient descent for non‑convex $f$ over a convex feasible set.
\item \textbf{L1‑regularization $g(x)=\lambda\|x\|_{1}$.} The proximal operator is soft‑thresholding; the bound (3) provides a guarantee for sparse non‑convex learning problems.
\end{enumerate}

\end{document}